% Generated draft: requires mathematical and source review.
\documentclass[11pt]{article}
\usepackage[margin=25mm]{geometry}
\usepackage{fontspec}
\setmainfont{Latin Modern Roman}
\usepackage{amsmath,amssymb,mathtools,unicode-math}
\setmathfont{Latin Modern Math}
\usepackage{enumitem,needspace,xurl,hyperref}
\setlength{\parindent}{0pt}
\setlength{\parskip}{5pt}
\setlength{\emergencystretch}{3em}
\newcommand{\N}{\mathbb{N}}
\newcommand{\Z}{\mathbb{Z}}
\newcommand{\Q}{\mathbb{Q}}
\newcommand{\R}{\mathbb{R}}
\newcommand{\C}{\mathbb{C}}
\newcommand{\F}{\mathbb{F}}
\newcommand{\A}{\mathbb{A}}
\newcommand{\PP}{\mathbb{P}}
\newcommand{\E}{\mathbb{E}}
\newcommand{\eps}{\varepsilon}
\newcommand{\dd}{\,\mathrm{d}}
\begin{document}
\section*{UnsolvedMath ID: 30002012}
\textbf{Source problem code:} catalogue entry
\subsection{Definitions and mathematical statement}
The supplied record concerns dense granular materials. Here \emph{dilatancy} and \emph{compaction} refer descriptively to increases and decreases in bulk volume during deformation. The source record does not specify a mathematical definition of the \emph{critical state}, a criterion for a \emph{persistent shear band}, or a definition of the \emph{intrinsic length scale} associated with such a band.

The recorded research aims are:
\begin{enumerate}
\item Explain the micromechanisms of dilatancy and compaction in dense granular materials.
\item Determine why the critical state depends on pressure.
\item Determine the intrinsic length scale governing persistent shear bands.
\end{enumerate}
These aims supply neither a constitutive or microscopic model nor material parameters, loading conditions, boundary conditions, a pressure regime, or quantified mathematical assertions. The pressure dependence is part of the recorded question; its applicability is not established here. No universal pressure law or length-scale formula can be extracted from the supplied information.

\textbf{Unresolved status of the intended statement.} The registry labels this entry \emph{open}, but also marks its statement \emph{unrecoverable}. Its supplied statement-verification note says that the official contribution introduces these physical questions and then reports results explaining the mechanisms, pressure dependence, and length scale. According to that note, the extraction summarizes answered research aims rather than a separate open problem. Consequently, the aims above preserve the source record but do not constitute a recovered, precisely quantified open mathematical problem.
\subsection{Short English statement}
The record asks for explanations of volume changes, pressure-dependent critical states, and the length scale of persistent shear bands in dense granular materials. However, its supplied verification note identifies these as research aims addressed by the contribution. A separate open problem and the hypotheses needed for a precise mathematical formulation have not been recovered.
\subsection{Sources}
\begin{itemize}
\item Supplied local UnsolvedMath registry record, permanent ID \(30002012\), title \emph{Dense-Granular-Flow Results Report}, exact source code \texttt{catalogue entry}, in the Oberwolfach Reports Open Problems collection. The record gives the proposed year as 2012 and names no individual proposer. It retains the original record publicly for provenance and promotes no separate clean statement.
\item \emph{Mechanics of Materials} (2012), Oberwolfach Reports, pp.~885--961, DOI 10.4171/owr/2012/15. Supplied source report and literature reference labelled ``OWR report'': \url{https://doi.org/10.4171/owr/2012/15}. The supplied background describes Oberwolfach Reports as documenting mathematical workshops and their open-problem discussions. No particular contribution, author, or problem-specific page is identified in this record.
\item Supplied literature reference: \emph{Dynamic length scales in crystal plasticity}, \url{https://arxiv.org/abs/2406.09892}. Its applicability to the granular-material questions is not established by the supplied record.
\item Literature assessment dated exactly \texttt{2026-08-21}: the supplied curation labels the topic open and describes critical-state pressure dependence and the governing intrinsic length as model- and material-dependent research questions. It states that current dislocation-continuum work supplies mechanisms and scales but no general determination. This is supplied dated curation derived from the OWR record, not an independently verified literature assessment; its time-sensitive status and cited claims require independent verification. It must be read alongside the conflicting supplied statement-verification note described above.
\end{itemize}
\end{document}
